\subsection{両純top零のH3領域にはJordan部分がない}
\label{subsec:h3-order2-22-rank4-zero-top-unipotent}
\begin{theorem}\label{thm:h3-order2-22-rank4-zero-top-no-Jordan}
$(2,2)$・奇関係数4のH3包絡$B$について、
$H=\mu_{B,u}^{1/4}$、$C=B^H$と置く。
半単純包絡$C$の純正・純負topがともに零なら、$H=I$である。
従ってこのcanonical零top領域の全H3は半単純中山を持つ。
全AS4の同領域に単冪自己同型がないという主張ではない。
\end{theorem}
\begin{proof}
$C$は半単純中山$(+,+,-,-)$を持ち、$H$は$C$の単冪自己同型でもある。
前節までに$C$の奇$UV$主block $A$は可逆交代行列である。
$H$のtop直線への作用は単冪、従って1なので、$H_U\otimes H_V$は$A$を保存する。

$H_U,H_V$がともに非自明なら、それぞれを
$J=\left(\begin{smallmatrix}1&1\\0&1\end{smallmatrix}\right)$にできる。
一般六係数交代行列について
\[
 (J\otimes J)A(J\otimes J)^t=A
\]
を解くと、辞書順六係数は$(z_0,z_1,0,0,0,0)$に限られ、Pfaffianは零となる。
これは$A$の可逆性に反する。
従って$H$が非自明なら片側だけが非自明である。

$C$の関係は二つの純AS2二次式と可逆mixed graphである。
逆twistは$\mathcal R_B=(1\otimes H)\mathcal R_C$を与えるので、
元の$B$にも二つの非退化純AS2関係と可逆graphが残る。
特に$(U,V)$二重次数があり、AS4の三次次元20とmonic簡約から
二つのAS2標準語を正・負の順に並べた全次数PBW基底を得る。
以下では$G$を$C$へ移さず、元の$B$の中山と整域性を使う。

まず$H_U\ne I,H_V=I$の場合を考える。
$\mu_B|_U$は非自明単冪、$\mu_B|_V=-I$である。
非自明単冪を持つ二元AS2平面は可換またはJordanであり、
$x$を固定し$y\mapsto y+\eta x$、$\eta\ne0$とする形にできる。
その次数$n$固定部分は$kx^n$である。
$\mu_B(G)=G$を正三次のPBW表示
\[
 G=H_3(U)+\sum_{w\in\{u^2,uv,v^2\}}a_w(U)w
\]
へ適用すると、$H_3\in kx^3$、全$a_w\in kx$となる。
従って$G$は$x$を一次因子に持ち、$B/(G)$は整域にならない。

次に$H_U=I,H_V\ne I$とする。
$H_V(u)=u,H_V(v)=v+u$と選べる。
可逆mixed graphがこのshearと可換する条件から
\[
 ut=f(t)u,\qquad vt=f(t)v+d(t)u\quad(t\in U)
\]
となる。実際、負次数を二blockに並べたgraphの可換環は
$\left(\begin{smallmatrix}f&d\\0&f\end{smallmatrix}\right)$である。
純負関係はshear固定tensorなので、非退化性から
$vu-uv=\kappa u^2$と書ける。
従って正AS2平面$T$の上の逐次Ore表示
\[
 S=T[u;f],\quad B=S[v;\sigma,\delta],\quad
 \sigma|_T=f,\quad\sigma(u)=u,\quad
 \delta(t)=d(t)u,\quad\delta(u)=\kappa u^2
\]
を得る。
ここで$f$の関係保存とOre条件は、PBW認識と三次結合則から従う。
元の中山は$\mu_B|_U=I,\mu_B(u)=-u$なので
\[
 f^2=\mu_T,\qquad\operatorname{hdet}_T(f)=-1
\]
となる。

$T$が可換なら$f$は反射である。
$d$は$f$と可換するので、$f(x)=x,f(y)=-y$の基底で$d$は対角である。
よって$y$は正規で、$B/(y)$は可換$k[x,u]$の通常微分Ore整域となる。
必要な$Gu=-uG$をこの商へ移し、$v$最高次数を比較すると$G$の像は零である。
従って$G$は$y$で割れ、整域商を与えない。

$T$が非可換量子平面なら、$G$の$V$次数2の最高$v$次数係数$z\in T_1$は
$z\mu_T(t)=tz$を全$t\in T_1$について満たす。
そのような一次式は零なので、順に全$V$次数2部分が消える。
残る$G\in Z(T)_3$は、非零なら二元AS2の一次因子分解により積となる。
$T$がJordanなら、$f^2=\mu_T$の半単純部分から$\operatorname{hdet}f=1$となり、
負中山に必要な$-1$と矛盾する。
全場合で非自明片側Jordanも排除され、$H=I$が従う。
\end{proof}

検証器は
\[
 \texttt{verify\_cubic\_order2\_22\_rank4\_zero\_top\_unipotent.py}
\]
である。両側単冪でのPfaffian消滅、片側graphの可換環、純負固定二次式を再生する。
AS2 shearの次数3固定核と一次intertwinerは既存のseparated検証器でも全係数確認済みである。
右twistの指数は$-4$であり、twistがH3を保つという仮定は使っていない。
